\section{Complete-trace equivalences coincide }
\label{s:wbc_vs_c}

In this section, we show that complete-trace equivalences 
in \COGS, \AOGS, and \WBOGS  coincide.

\begin{remark}
Complete trace equivalence (in all the OGS forms discussed)  is strictly weaker than
ordinary trace equivalence, as shown by the following terms:
\[
 (\lambda z.\Omega)(x (\lambda y.\Omega)) \text{ and }
 (\lambda z.\Omega)(x (\lambda y.y))
\]
The terms  are equal  if only complete traces are observed, while  they are
not otherwise.
\end{remark}


\reviewComment{
 I am confused here. This example should be more
detailed, because to me these configurations do *not* seem to be trace
equivalent in WB-OGS. There seems to be a tension with Lemma 10.4.
}
\reviewCommentDS{I have rephrased the remark, see if ok; indeed i believe the example was about
the difference between complete and ordinary trace equivalence, something i guess well known}
\begin{theorem}
\label{t:main_2_WBC}
For any $\FFb,\GGb$, 
we have 
$\FFb  \TEwbc \GGb$
iff
$\es\FFb  \TEc \es\GGb$ iff
$\es\FFb  \TEpc \es\GGb$.
\end{theorem} 
\noindent
One direction of the theorem is easy, intuitively because a trace in 
 \WBOGS is 
also a trace in \AOGS, and these, in turn,  are traces in 
 \COGS.
The problem is the converse, namely showing 
that  
$\FFb  \TEwbc \GGb$
implies
$\es\FFb  \TEc \es\GGb$, and 
that 
$\es\FFb  \TEc \es\GGb$ 
implies $\es\FFb  \TEpc \es\GGb$.
They are proven in the next two sections.

\begin{remark}
\label{r:CT_COGS}
In earlier sections, (Section~\ref{s:cogs}),
we have  derived another equality result for
trace equivalence, namely Corollary~\ref{c:OGS_COGS}, 
relating  trace equivalences in the Alternating 
and Concurrent LTS. 
The structure of the proofs of that result and the one above 
(Theorem~\ref{t:main_2_WBC})
 are however quite different. 
In  Corollary~\ref{c:OGS_COGS} we derived the result via up-to techniques; 
we cannot apply such techniques here because we are reasoning on trace equivalence,
notably a restricted form of traces (the complete ones). 
On the other hand,  the proof strategy here 
rests on decomposition properties 
such as  Lemma \ref{l:decWBCN} that rely on a termination property for initial
configurations, which may not be assumed in 
Section~\ref{s:cogs}.
\end{remark} 

\subsection{Relating \WBOGS and \AOGS complete traces}
\label{ss:main_wbct}

\ifcutLMCS


\begin{lemma}
\label{l:FFbFF}
For  configurations $\stack \FFb \stt$ and $\stack \GGb \stt$, if 
$ \es\FFb  \TEcN  \stt  \es\GGb$,  then $
\FFb  \TEwbcN \stt \GGb$. 
\end{lemma} 

The problem is the converse. 
This is tackled in the remainder of the section 
(the main result is Lemma~\ref{l:mainWBCN}).
\fi
             
\begin{lemma}
\label{l:TCtoWB}
Taking $\FFb$ a \WBOGS configuration,
if $\es\FFb$ has a complete trace in \AOGS, 
then $\FFb$ has a well-bracketed complete trace. 
\end{lemma} 


\begin{proof}
Assume $\trace$ is a complete trace for $\es\FFb$ in \AOGS.
We proceed by induction on the length of $\trace$. 
If $\trace$ is empty then the set of P- and O- continuation names of
$\es\FFb$ is empty.
So $\trace$ is also a complete trace for $\FFb$
since the stack of $\FFb$ must then be empty.

Suppose now $\trace = \mu \trace'$ is not empty,
so that there exists $\FF'$
such that $\es{\FFb} \ArrA \mu \FF' \ArrA{\trace'}$.
We reason by case analysis over $\mu$:
\begin{itemize}
\item If $\mu$ is a question,  then $\mu$ is either 
of the form $\questPV{x}{y}{p}$ 
or of the form $ \questOV{x}{y}{p}$.
We also have $\FFb  \ArrWB \mu \FFb'$
with 
$\es{\FFb'}= \FF'$.
Moreover $\trace $ is a complete trace %\wrt  $p::\stt$ 
for $\FF'$. 
By induction $\FFb'$ 
has a well-bracketed complete trace, %\wrt $p::\stt$, 
say $\tracerho$.
Then $\mu \tracerho$ is a well-bracketed complete trace %\wrt $\stt$ 
for $\FFb$. 
\item Suppose now $\mu$ is an answer. There are two cases:
\begin{itemize}
\item If $\FFb$ is active, then the answer $\mu$ must be of the form 
$\ansP{p}{x}$, where $p$ is the continuation name for the running term.
%, and  $\stt = p::\sttp$. 
Thus $\mu$ is also the correct answer for \WBOGS, and we can easily conclude using
induction. 
\item If $\FFb$ is passive, then 
 $\mu$ must be of the form 
 $\ansO{p}{x}$. Writing $\stt$ for the stack of $\FFb$,
 if $\stt = p::\sttp$ then $\mu$ is also    
 the correct answer for a well-bracketed trace, and as before  we can conclude using
 induction.  Otherwise, suppose 
 $\stt = q::\sttp$ with $q\neq p$
 (i.e., $p$ is not the last  name in $\stt$).
 This means that $\FFb$ has at least two
 singleton components, located at $p$ and $q$. 
 Take a decomposition
 $\FFb = \FFb_1 \otimes_{\sigma} \FFb_2 $,    
 where $p$ and $q$ end up in different sub-configurations. 
% Writing $\stti$ for the set of continuation names of $\FF_i$,
 We get from Lemma~\ref{l:AOGS-decomp-ct}
 that there exists $\trace_i$ complete for $\es{\FFb_i}$ in \AOGS, 
 %\wrt $\stti$ with $\es{\FFb_i} \ArrA {\trace_i}$, 
 for $i\in \{1,2\}$, and 
 $\mu \trace \in \interlA{\trace_1}{\trace_2}$.
 By induction, there is $ \tracerho_i $ a complete well-bracketed trace 
 for $\FFb_i$ in \WBOGS.
%\wrt $\stti$,  such that 
%$\FFb_i \ArrWB {\tracerho_i} $.
%
 By Lemma~\ref{c:dec_GS_WB}, 
 there is $\tracerho \in \interlWBs{\stt}{\tracerho_1}{\tracerho_2}$, 
 a complete well-bracketed trace for $\FFb$.
% \wrt $\stt$,
% such that $\FFb \ArrWB \tracerho$.
\qedhere
\end{itemize}
\end{itemize}
\end{proof} 

\begin{lemma}
\label{l:mainWBCN}
Suppose 
$\FFb \TEwbc \GGb$,
then
$\es\FFb \TEc \es\GGb$
\end{lemma}              

\begin{proof}
First, if $\FFb,\GGb$ are not terminating, i.e. they have no complete trace in \WBOGS, 
then from Lemma~\ref{l:TCtoWB}, both 
$\es{\FFb},\es{\GGb}$ have no well-bracketed complete trace, thus
$\es\FFb \TEc \es\GGb$. 

Suppose now that $\FFb$ and $\GGb$ are both terminating.
By Lemmas~\ref{l:decWBCN} and \ref{l:compCN}, it is sufficient to consider 
the case when $\FFb$ and $\GGb$ are singleton. 
%
Suppose there exists a complete trace $\tr$ for $\es{\FFb}$, 
then we prove by induction over $\tr$
that $\tr$ is a complete trace for $\es{\GGb}$ as well.
If $\tr$ is empty, then from Lemma~\ref{l:OGS_complete-strongly-passive},
$\es{\FFb}$ must be strongly passive.
Recall that from $\FFb \TEwbc \GGb$ we enforce that
$\okC{\FFb}{\GGb}$.
So $\es{\GGb}$ must be strongly passive as well, so that 
the empty trace is complete for $\es{\GGb}$ as well.

Suppose now that 
$\tr = \act \tr'$, so that
$\es\FFb \Longrightarrow \arrA {\act} \FF' \ArrA{\tr'} \FF''$
with $\tr'$ complete for $\FF'$.
Then $\act$ is the only possible visible transition for $\es\FFb$ 
since $\es\FFb$ is a singleton.
So it is also the first
visible transition in any well-bracketed complete trace for $\FFb$ (i.e., we have 
$\FFb \ArrWB{} \arrWB{\act} \FFb'$ with $\es{\FFb'} = \FF'$). 
Indeed, if $\act$ is an Opponent answer $\ansO{p}{x}$, then the stack
of $\FFb$ must be equal to $p$ since $\FFb$ is a singleton.
Since $\FFb$ is terminating, it has thus a well-bracketed complete
trace starting with $\act$.
%
So does $\GGb$ since 
$\FFb  \TEwbc \GGb$, that is, we have  
$\GGb \ArrWB{} \arrWB{\act} \GGb'$, for some $\GGb'$, and therefore also 
$\es\GGb \ArrA{} \arrA{\act} \es{\GGb'}$. 
% 
Moreover, $\FFb'  \TEwbc \GGb'$.
% 
% , where $\sttp$ is obtained from $\stt$ and $\mu$ in
% the expected manner.
% \finish{this could be shown as a separate lemma, but it seems straightforward, as the
%   first action $\mu$ is deterministic}  

If $\FFb'$ and $\GGb'$ are singleton, we can now apply induction and 
derive that $\tr'$ is also complete for $\es{\GGb'}$, and then we are done. 
%
If they are not singleton, we can consider a decomposition for them 
(yielding $\FFb'_i$ and $\GGb'_i$), and associated decomposition
for $\tr'$ in complete traces $\tr'_i$ for $\es{\FFb_i}$ using Lemma~\ref{l:AOGS-decomp-ct}. 
We can then apply induction on the singleton components (using
Lemma~\ref{l:decWBCN}), then we infer that 
$\tr'_i$ is also a complete trace for $\es{\GGb_i'}$. 
We can now apply Lemma~\ref{c:dec_GS} and infer that $\tr'$ is a complete trace for $\es{\GGb'}$,
and then conclude that $\act\tr'$ is a complete trace for $\es\GGb$.
\end{proof} 

\subsection{Relating \AOGS and \COGS complete traces}
\label{a:aogs-cogs-complete}


\begin{lemma}
\label{l:CtoCA}
Taking $\FF$ an \AOGS configuration,
if $\FF$ has a complete trace in \COGS, then $\FF$ 
has a complete trace in \AOGS. 
\end{lemma}

\begin{proof}
Assume $\trace $ is a complete trace
for $\FF$ in \COGS.
We proceed by induction on the length of $\trace $.
If $\trace$ is empty then this is straightforward.
Suppose now there exists $\mu$ s.t. 
$\trace = \mu \tr'$. We reason by case analysis on $\mu$:
\begin{itemize}
\item If $\mu$ is a Player action,
then $\FF \ArrA \mu \FF'$, so we conclude using induction.
\item If $\mu$ is an Opponent action,
there are two cases:
\begin{itemize}
\item If $\FF$ is a passive configuration, 
then again $\FF \ArrA \mu \FF'$, and we conclude using induction.
\item If $\FF$ is an active configuration, then we can write
$\FF$ as $\FF_1 \otimes \FF_2$ since $\FF$ has both
an active term and at least one element in its environment,
that is triggered by $\mu$.
%Writing $\Sigma_i$ for the set of continuation names of $\FF_i$,
We get from Lemma~\ref{l:COGS-decomp-ct}
that there exist two traces $\trace_1,\trace_2$ complete respectively for
$\FF_1$ and $\FF_2$ in \COGS such that $\mu \trace \in \interl{\trace_1}{\trace_2}$.
By induction, there exists two traces $\tracerho_1,\tracerho_1$ 
complete respectively for $\FF_1$ and $\FF_2$ in \AOGS.
%
Finally,
by Lemma~\ref{l:AOGS-recomp-ct}, there exists a trace
$\tracerho \in \interlA{\tracerho_1}{\tracerho_2}$
that is complete for $\FF$ in \AOGS.
\qedhere
\end{itemize}
\end{itemize}
\end{proof} 

\begin{lemma}
\label{l:mainCC-CA}
Suppose 
$\FF  \TEc \GG$,
then
$\FF  \TEpc \GG$.
\end{lemma}   

\begin{proof}
First, if $\FF,\GG$ are not terminating, i.e. they have no complete trace in \AOGS, 
then from Lemma~\ref{l:CtoCA}, both 
$\FF,\GG$ have no complete trace \COGS, thus
$\FF \TEpc \GG$. 

Suppose now that $\FF$ and $\GG$ are both terminating in \AOGS.
By Lemmas~\ref{l:decACN} and the adaptation of Lemma~\ref{l:compCN} to \COGS, it is sufficient to consider 
the case where $\FF$ and $\GG$ are singleton. 
%
Suppose there exists a complete trace $\tr$ for $\FF$ in \COGS, 
then we prove by induction over $\tr$
that $\tr$ is a complete trace for $\GG$ in \COGS as well.
If $\tr$ is empty, then $\FF$ must be strongly passive.
Recall that from $\FF \TEc \GG$ we enforce that
$\okC{\FF}{\GG}$.
So $\GG$ must also be strongly passive, so that 
the empty trace is a complete trace for $\GG$ as well.


If $\tr$ is of the shape $\act\tr'$, then
$\FF \Longrightarrow \arrC {\mu} \FF' \ArrC{\tr'}$.
Since $\FF$ is singleton, $\FF$ cannot be active with $\mu$ an Opponent action.
So $\mu$ is a Player action, or an Opponent action with $\FF$ passive,
so that $\FF \ArrA{}  \arrA {\mu} \FF'$.
Since $\FF$ is terminating in \AOGS, we get that $\FF$ has a complete trace
starting with $\mu$ in \AOGS.
Since 
$\FF  \TEc \GG$, we also have
$\GG \ArrA{} \arrA \mu \GG'$, for some $\GG'$, and therefore also 
$\GG \ArrC{} \arrC \mu \GG'$. 
Moreover, $\FF' \TEc \GG'$.

If $\FF'$ and $\GG'$ are singleton, we can now apply the induction hypothesis 
and derive the same complete trace  $\tr'$ also for $\GG'$ in \COGS, and then we are done. 
%
If they are not singleton, we can consider a decomposition for them 
(yielding $\FF'_i$ and $\GG'_i$), and associated decomposition
for $\tr'$ in complete traces $\tr'_i$ for $\FF_i$ using Lemma~\ref{l:COGS-decomp-ct}. 
We can then apply induction on the singleton components (using
Lemma~\ref{l:decACN}), and getting that
$\tr'_i$ is also a trace for $\GG_i'$ in \COGS. 
We can now apply Lemma~\ref{l:COGS-recomp-ct} 
and infer that $\tr'$ is a complete trace for $\GG'$ in \COGS,
and then conclude that $\act\tr'$ is a complete trace for $\GG$ in \COGS.
\qedhere
\end{proof} 

